\documentclass{article}
\usepackage[dutch]{MJPreamble}

\title{Midterm Organisations and Markets 2026}
\author{}
\date{}

\begin{document}
\maketitle

\section{[70 points] Entry and Competition}
In a market for potatoes, demand is given by $D(p)=630-\frac p2$. If a firm produces potatoes, it has to pay a cost $c(q)=4q^{2}+100$ to produce $q$ potatoes, where the 100 is a fixed cost, only payable if the firm enters into the market.

\begin{enumerate}[label=(\alph*)]
\item \textit{9 points} Suppose there is a single price-taking firm. Calculate the equilibrium price and quantity. \end{enumerate}
\textbf{Solution}:
Inverse demand is $p=1260-2q$. \textit{(2 points)}
$\Pi(q)=pq-c(q)=pq-4q^{2}-100$, which gives first order conditions: $p-8q=0$ (if you find these from $p=mc$, that is fine too) \textit{(correct set-up: 4 points)}. Setting (inverse) supply and demand equal: $8q=1260-2q$, or $q=126$ and hence $p=1008$. \textit{(Correct calculation: 3 points)} $\square$

\begin{enumerate}[resume*]
\item \textit{4 points }What are the Lerner index and the Herfindahl-Hirschman index (measured from 0 to 1) for this equilibrium?\end{enumerate}
\textbf{Solution}:
$p=mc$, so the Lerner index is 0, as it always is in a price-taking market. \textit{(2 points)} The H-index is just 1, since there is only one firm. \textit{(2 points)} $\square$

\begin{enumerate}[resume*]
\item \textit{9 points }Now suppose that firms can enter at will, but are still price taking and have to pay the fixed cost for entering. Calculate the equilibrium price and quantities (per firm and in total).\end{enumerate}
\textbf{Solution}:
Write $q$ for the production of a single firm and $Q$ for the total production. If firms can enter, this drives profits to 0, or, equivalently, $mc=ac$. \textit{(4 points)} We already know from the first question that each firm sets $q$ so that $p=8q$. $ac=4q+\frac{100}q$. \textit{(2 points)} Equating them gives $4q^{2}=100$, or $q=5$. Then $p=40$ and $Q=630-\frac {40}2=610$. \textit{(3 points)}$\square$

\begin{enumerate}[resume*]
\item \textit{4 points }What are the Lerner index and the Herfindahl-Hirschman index (measured from 0 to 1) for this equilibrium?\end{enumerate}
\textbf{Solution}:
$p=mc$ still, so the Lerner index remain 0. \textit{(1 point)} Each firm produces $q=5$ and the total produce is $610$, so there are $n=610/5=122$ firms. Since they are all equally sized, $H=\frac 1n=\frac 1{122}$. \textit{(3 points)} $\square$

\begin{enumerate}[resume*]
  \item \textit{9 points }How do the two equilibria compare to each other for the consumers? For the firms? In terms of efficiency (welfare)?\end{enumerate}
\textbf{Solution}:
The consumers like the second better, as prices are lower (and quantities sold higher). \textit{(3 points)} Firms like it less, as their profit is driven to 0 \textit{(2 points)}. Both situations are efficient with respect to their productive capacity, but in the second there are more firms, so the productive capabilities are improved. \textit{(4 points)} $\square$


From now on assume firms are not price-taking, but strategically competing by setting quantities.

\begin{enumerate}[resume*]
\item \textit{9 points }Suppose again that there is a single firm. Calculate the equilibrium price and quantity.\end{enumerate}
\textbf{Solution}:
A monopolist exploits his understanding of the demand function. Recall that $p=1260-2q$. \textit{(2 points if not earned in 1.a)}
$\Pi(q)=p(q)q-c(q)=(1260-2q)q-4q^{2}-100$, \textit{(4 points)} which gives first order conditions $1260-4q-8q=0$ (if you find these from $mr=mc$, that is fine too). Thus $q=105$ and hence $p=1050$. \textit{(5 points)} $\square$

\begin{enumerate}[resume*]
\item \textit{4 points }What are the Lerner index and the Herfindahl-Hirschman index (measured from 0 to 1) for this equilibrium?\end{enumerate}
\textbf{Solution}:
The marginal cost is $mc=8q=840$. $L=\frac{p-mc}p=\frac{210}{1050}=0.2$. \textit{(3 points)} There is only one firm, so the H-index is 1. \textit{(1 point)} $\square$

\begin{enumerate}[resume*]
\item \textit{9 points }Now suppose there are two firms that simultaneously set a quantity to produce. Calculate the equilibrium price and quantity.\end{enumerate}
\textbf{Solution}:
Each firm has profit function $\Pi_{i}(q_{i};q_{\neg i})=p(q_{i},q_{\neg i})q_{i}-c(q_{i})=(1260-2q_{i}-2q_{\neg i})q_{i}-4q_{i}^{2}-100$. \textit{(4 points)} This gives first order conditions: $1260-2q_{\neg i}-4q_{i}-8q_{i}=0$. \textit{(2 points)} Since the model is symmetric, we know $q_{i}=q_{\neg i}$, so $q_{i}=\frac{1260}{14}=90=q_{\neg i}$. So the price is $p=1260-2Q=900$. \textit{(3 points)} $\square$

\begin{enumerate}[resume*]
\item \textit{3 points }What is the Herfindahl-Hirschman index (measured from 0 to 1) for this equilibrium?\end{enumerate}
\textbf{Solution}:
There are two equal sized firms, so $H=\frac 12$. $\square$

Now suppose that firms can enter into this market at will (while still competing in quantity and still facing a fixed cost of 100).

\begin{enumerate}[resume*]
  \item \textit{Bonus: 10 points }Suppose in this new situation the market demand increases to $D_{New}(p)=1260-p$. This will have two distinct effects on efficiency. Explain intuitively what they are and what their direction is. Show no calculations. 
  \end{enumerate}
  \textbf{Solution:}
  Increasing the market size will attract more firms. One the one hand, this will increase competition, leading to lower prices. Since $p>mc$, this increases efficiency. \textit{(5 points)} On the other hand, every firm entering has to pay the fixed cost, which is wasteful. \textit{(5 points)} $\square$
  
\section{[40 points] Teams}

Alice and Bob can work together to produce widgets, which they can sell at price 1. If Alice works $A$ hours and Bob works $B$ hours, they produce $Q(A,B)=30A+30B+AB$ widgets. Their cost of working is $C(A)=A^{2}$ and $C(B)=B^{2}$ respectively. 

\begin{enumerate}[label=(\alph*)]
\item \textit{8 points }Compute Alice's marginal productivity. What can you say about how Alice and Bob produce as a team? \end{enumerate}
\textbf{Solution}:
$MP_{A}=\frac{\partial Q}{\partial A}=30+B$. \textit{(4 points)} Alice's marginal productivity increases when Bob works more. They make a good team. \textit{(4 points)} $\square$

\begin{enumerate}[resume*]
\item \textit{8 points }What would be the efficient levels of $A$ and $B$? And the corresponding level of output?\end{enumerate}
\textbf{Solution}:
$V(A,B)=30A+30B+AB-A^{2}-B^{2}$, \textit{(3 points)} which gives first order conditions $30+B-2B=0$ and $30+A-2B=0$ \textit{(2 points)}, so $A=B=30$ and total produce is $2700$. \textit{(3 points)} $\square$

\begin{enumerate}[resume*]
\item \textit{8 points }Suppose that Alice and Bob split the revenue equally and simultaneously choose how many hours to work. What will be the resulting equilibrium hours worked and equilibrium output?\end{enumerate}
\textbf{Solution}:
Alice has pay-off $\Pi_{A}(A;B)=\frac{30A+30B+AB}2-A^{2}$, \textit{(3 points)} which gives first order condition: $15+\frac B2-2A=0$. \textit{(2 points)} A similar calculation for Bob (or symmetry implying $A=B$) gives $A=B=10$ and total produce is $700$. \textit{(3 points)} $\square$

Now suppose Chrysostomos offers a prize $P$ to Alice and Bob. Only one of them can get the prize and the probability that Alice will get it is $\dfrac{A}{A+B}$ and the probability that Bob will get it is $\dfrac B{A+B}$. Chrysostomos gets to keep and sell whatever Alice and Bob produce.

\begin{enumerate}[resume*]
  \item \textit{8 points }What does Alice's pay-off function look like in this scenario? And Bob's pay-off function?\end{enumerate}
  \textbf{Solution}:
  $P_{A}=P\frac A{A+B} -A^{2}$. $\square$

\begin{enumerate}[resume*]
  \item \textit{8 points }How high does Chrysostomos have to set $P$ in order to induce the efficient production level?\end{enumerate}
\textbf{Solution}:
From the previous question, we get Alice's first order condition: $P\frac B{(A+B)^{2}}-2A=0$. \textit{(2 points)} By symmetry, $A=B$ and $P=8A^{2}$. \textit{(3 points)} Since Chysostomos wants to induce $A=30$, he has to set $P=7200$. \textit{(3 points)} (That is substantially more than the value of the produce, which is $2700$. \textit{(2 points bonus)})

\end{document}
